\section{Proof audit and technical checkpoints}\label{app:audit}

This appendix records the points at which a formal variational or finite--$N$ calculation can fail and explains how they are handled in the proof above.  It is included deliberately: the exact numerical value of $\beta_3$ is meaningful only if the radial variational argument is genuinely global, and the improved atomic coefficient is meaningful only if the finite--$N$ propagation is numerically safe.

\subsection{Enlarging the HPS admissible class}

The HPS definition uses radial probability measures belonging to $H^{-1}(\R^3)$ and with finite second moment.  For the lower bound we enlarged the class after radial reduction.  This is harmless because the lower bound is proved for the larger class, while the explicit optimizer \eqref{eq:rho-star} is an $H^{-1}$ probability density with compact support bounded away from the origin.  Hence the lower and upper bounds meet inside the original HPS class.  HPS radial reduction identifies the value of the full variational infimum with the radial infimum; the uniqueness argument in the present manuscript, however, is only a uniqueness statement inside the radial class.

\subsection{Existence and loss of exponential moments}

Weak compactness alone does not preserve the constraints
\[
 A[\sigma]=B[\sigma]=1,
\]
because exponentially small masses can escape to large $|t|$ while carrying a nonvanishing exponential moment.  The proof of \Cref{lem:existence} therefore uses more than tightness.  First, a normalized two--point measure gives an explicit competitor with $J>0$, so the supremum $M$ is strictly positive.  Tightness and bounded total mass yield narrow convergence along a maximizing subsequence, and the corresponding product measures converge narrowly as well; because $k(t-u)$ is bounded and continuous, this gives $J[\sigma_n]\to J[\sigma]=M$.  If the weak limit had $A B<1$, amplitude and translation normalization would multiply $J$ by $(AB)^{-1}>1$, producing a normalized competitor with value larger than $M$.  This closes the only possible exponential-moment loss mechanism.

\subsection{The obstacle condition without an abstract KKT theorem}

The first variation is derived from the scale-- and translation--invariant quotient
\[
 \frac{J[\sigma]}{A[\sigma]B[\sigma]}.
\]
Adding an atom $\varepsilon\delta_t$ gives the global inequality directly.  Equality on the support then follows from the identity
\[
 \int\left[M(\ee^t+\ee^{-t})-2(k*\sigma)(t)\right]\dd\sigma(t)=0
\]
and nonnegativity of the bracket.  Thus no constraint qualification for an infinite-dimensional KKT theorem is required.

\subsection{Atoms and smooth pasting}

The cusp of $k$ at the origin is essential.  An atom of mass $m$ in $\sigma$ produces an upward derivative jump $2m$ in $k*\sigma$ and hence a downward jump $-4m$ in the derivative of the nonnegative obstacle.  This is incompatible with contact at a local minimum, excluding atoms.

For a compactly supported atomless finite measure, convolution with the Lipschitz kernel $k$ is $C^1$: the one-sided derivative discontinuity of $k'$ occurs only at $u=t$, which carries zero mass, and continuity follows by dominated convergence after splitting off a small neighborhood of $t$.  Therefore the obstacle is $C^1$, and at either extremal contact point one has both value and derivative equal to zero.  This justifies the smooth-pasting equations used in \Cref{lem:endpoints}.

\subsection{Conditional negativity for signed measures}

For a finite signed measure $\eta$ satisfying the two moment constraints, the function
\[
 p(t)=\int\ee^{|t-u|}\dd\eta(u)
\]
vanishes outside the candidate interval and belongs to $H_0^1$.  Distributionally,
\[
 (\partial_t^2-1)p=2\eta,
\]
so in particular $\eta\in H^{-1}(\R)$.  Fourier transformation gives
\[
 \widehat\eta(\xi)=-\frac{1+\xi^2}{2}\widehat p(\xi).
\]
Combining this identity with
\[
 \widehat{k}(\xi)=\frac{2(2-\xi^2)}{(1+\xi^2)(4+\xi^2)}
\]
gives the multiplier decomposition
\[
 \widehat{k}(\xi)|\widehat\eta(\xi)|^2
 =\frac12\left(-\xi^2+5-\frac{18}{\xi^2+4}\right)|\widehat p(\xi)|^2,
\]
which is exactly \eqref{eq:Jeta-identity}.  The finite-measure pairing can be justified by standard $H^{-1}$ approximation.  Thus the signed-measure step does not rely on a formal Fourier transform of a measure without a function-space interpretation.  The strict inequality ultimately uses only
\[
 2a_*<\frac{\pi}{\sqrt5}.
\]

\subsection{Why the global argument does not assume connected support}

No connectedness assumption on $\supp\sigma$ is made.  We define only its extremal points $a$ and $b$.  Compactness and smooth pasting at these two points yield the endpoint identities and hence $a+b=0$.  If a maximizer had value at least $M_*$, its extremal support points would lie inside $[-a_*,a_*]$, regardless of holes in the support.  The difference measure $\eta=\sigma-\sigma_*$ is therefore supported in the full candidate interval, exactly the hypothesis required in \Cref{lem:conditional-negativity}.

\subsection{Scope of the uniqueness statement}

The proof of \Cref{thm:main} shows that the normalized maximizer of the logarithmic radial problem is unique, and hence that the radial minimizer is unique up to dilation.  HPS radialization at $s=3$ proves equality between the full and radial infima, but the radialization inequality is not used here in a strict form.  Consequently the present paper does not claim that every minimizer in the full nonradial HPS admissible class must itself be radial.

\subsection{Why the exact $\beta_3$ may be inserted into the HPS finite--$N$ argument}

The finite--$Z$ theorem is not obtained by replacing a printed decimal inside the final HPS statement.  What is used is the symbolic cubic-weight chain in \cite[Section 7.3]{HPS2025}.  In particular, their finite--configuration estimate has the form
\[
 (1-r^2)N\beta_3-\frac1r
 \le
 \left(1+\frac{r^2}{5}+\frac{r^3}{35}\right)\alpha_{N,3}(N-1),
\]
and their kinetic estimate bounds the right--hand side by the same factor times
\[
 Z+c_0ZN^{-2/3}.
\]
Thus the true mean--field constant $\beta_3$ remains symbolic throughout the finite--$N$ argument.  The old pointwise lower bound is used only to turn $\beta_3^{-1}$ into the decimal $1.1185$ and to produce safe decimal lower--order coefficients.  Replacing that lower estimate by the exact value from \Cref{thm:main} is therefore logically legitimate.

\subsection{Endpoint check in the $Z^{1/3}$ optimization}

After the lower--order coefficients are removed, the HPS contradiction argument reduces to bounding
\[
 F_{\beta_3}(x)
 =
 3\left(\frac3{10}\right)^{1/3}\beta_3^{-2/3}x^{1/3}
 +c_0\beta_3^{-1}x^{-2/3},
 \qquad x=\frac NZ.
\]
Because
\[
 F'_{\beta_3}(x)
 =\frac{x^{-5/3}}{3}(Ax-2B)
\]
with $A,B>0$, the function has a unique interior minimum and no interior maximum.  Hence one must compare both endpoints of the admissible interval.

For the present theorem $Z\ge4$.  Lieb's bound gives
\[
 x<2+\frac1Z\le\frac94,
\]
so the relevant interval is $[\beta_3^{-1},9/4]$.  Directed evaluation of the explicit expressions gives
\[
 F_{\beta_3}(\beta_3^{-1})
 =3.81147042506\ldots,
 \qquad
 F_{\beta_3}(9/4)
 =3.78401155302\ldots.
\]
Thus the left endpoint is the true maximum on the interval used here, and the rounded coefficient $3.812$ is safe.

This sharper interval is important for the numerical audit.  If one unnecessarily enlarges the interval to $[\beta_3^{-1},5/2]$, the right endpoint is slightly larger for the exact value of $\beta_3$.  The $Z\ge4$ theorem avoids that artificial loss by retaining the sharper consequence of Lieb's estimate.

\subsection{Certified decimal margins}

All decimals in \Cref{thm:finite-Z} are chosen strictly on the safe side of explicit closed expressions:
\[
 \beta_3^{-1}=1.08632517518\ldots<1.086326,
\]
\[
 \bar a_2=0.01293244256\ldots<0.01294,\qquad
 \bar a_3=0.18184084198\ldots<0.18185,
\]
\[
 \bar a_4=0.01940243839\ldots<0.01941,
 \qquad
 \sup F_{\beta_3}=3.81147042506\ldots<3.812.
\]
These are evaluations of elementary explicit formulas involving $\exp$, $\arctan$, radicals and $\pi$; no numerical optimization enters the theorem.

\subsection{What is not claimed}

The paper proves a sharpened finite--$Z$ theorem within the existing HPS cubic-weight framework, but it does not prove the ionization conjecture $N_c(Z)\le Z+C$.  Nor does it claim that the coefficient $\beta_3^{-1}$ is optimal among all possible weighted multiplier arguments.  Any further reduction of the leading coefficient must use information not present in the scalar radial mean--field problem, for example a quantitative stability term coupled to fermionic occupation constraints.
